\documentclass[aspectratio=169,10pt]{beamer}

% Shared Beamer setup for Fall 2026 course slides.
% Clean Cal Poly-inspired theme: no custom class files or uncommon packages.
\mode<presentation>{
  \usetheme{default}
  \usefonttheme{professionalfonts}
}

\definecolor{CalPolyGreen}{HTML}{154734}
\definecolor{CalPolyGold}{HTML}{BD8B13}
\definecolor{CalPolySage}{HTML}{7FA28F}
\definecolor{CalPolyMint}{HTML}{DDF1E7}
\definecolor{CalPolySand}{HTML}{A29F86}
\definecolor{CalPolyBeige}{HTML}{EFEEDE}
\definecolor{DeckInk}{HTML}{1F2933}
\definecolor{DeckMuted}{HTML}{5F6B7A}
\definecolor{DeckSoft}{HTML}{F7F7F5}

\setbeamersize{text margin left=0.55in,text margin right=0.55in}
\setbeamercolor{normal text}{fg=DeckInk,bg=white}
\setbeamercolor{structure}{fg=CalPolyGreen}
\setbeamercolor{title}{fg=white}
\setbeamercolor{frametitle}{fg=CalPolyGreen,bg=white}
\setbeamercolor{block title}{bg=CalPolySage,fg=white}
\setbeamercolor{block body}{bg=CalPolyMint,fg=black}
\setbeamercolor{alerted text}{fg=CalPolyGold}

\setbeamerfont{title}{size=\Huge,series=\bfseries}
\setbeamerfont{frametitle}{size=\Large,series=\bfseries}
\setbeamerfont{block title}{size=\normalsize,series=\bfseries}

\setbeamertemplate{navigation symbols}{}
\setbeamertemplate{itemize item}{\color{CalPolyGold}$\blacktriangleright$}
\setbeamertemplate{itemize subitem}{\color{CalPolyGreen}$\bullet$}
\setbeamertemplate{footline}{
  \hfill{\color{DeckMuted}\scriptsize\insertframenumber/\inserttotalframenumber}
  \hspace{0.18in}\vspace{0.12in}
}
\setbeamertemplate{frametitle}{
  \vspace{0.18in}
  \hspace{0.02in}\insertframetitle\par
  \vspace{0.08in}
  \noindent\makebox[\linewidth][l]{%
    {\color{CalPolyGold}\rule{0.55in}{2pt}}%
    {\color{CalPolyGreen}\rule{\dimexpr\linewidth-0.55in\relax}{2pt}}%
  }
}

\newcommand{\CourseNumber}{Course}
\newcommand{\CourseTitle}{Course Title}
\newcommand{\LectureNumber}{0}
\newcommand{\LectureTitle}{Lecture Title}
\newcommand{\LectureDate}{Date}
\newcommand{\CourseUnit}{Unit}

\newcommand{\coursenumber}[1]{\renewcommand{\CourseNumber}{#1}}
\newcommand{\coursetitle}[1]{\renewcommand{\CourseTitle}{#1}}
\newcommand{\lecturenumber}[1]{\renewcommand{\LectureNumber}{#1}}
\newcommand{\lecturetitle}[1]{\renewcommand{\LectureTitle}{#1}}
\newcommand{\lecturedate}[1]{\renewcommand{\LectureDate}{#1}}
\newcommand{\courseunit}[1]{\renewcommand{\CourseUnit}{#1}}

\author{Kirk Duran}
\institute{California Polytechnic State University, San Luis Obispo}
\date{\LectureDate}

\newcommand{\makedecktitle}{
  \begingroup
  \setbeamercolor{background canvas}{bg=CalPolyGreen}
  \begin{frame}[plain]
    \vfill
    \begin{center}
      {\usebeamerfont{title}\color{white}\LectureTitle\par}
      \vspace{0.18in}
      {\Large\color{white}Lecture \LectureNumber\par}
    \end{center}
    \vfill
    \noindent{\color{white}\rule{\linewidth}{1.2pt}}\par
    \vspace{0.08in}
    \noindent\colorbox{CalPolyGold}{\makebox[0.24\linewidth][c]{\color{white}\Large\CourseNumber}}
    \hspace{0.12in}{\color{white}\Large Kirk Duran}
  \end{frame}
  \endgroup
}

\newenvironment{keyidea}{\begin{block}{Key idea}}{\end{block}}
\newenvironment{activity}{\begin{block}{In class}}{\end{block}}
\newenvironment{cpdefinition}{\begin{block}{Definition}}{\end{block}}
\newenvironment{exampleblockcp}{
  \begingroup
  \setbeamercolor{block title}{bg=CalPolySand,fg=white}
  \setbeamercolor{block body}{bg=CalPolyBeige,fg=black}
  \begin{block}{Example}
}{
  \end{block}
  \endgroup
}

\newcommand{\imageplaceholder}[2][0.3in]{%
  \noindent\makebox[\linewidth][c]{%
    \fcolorbox{CalPolySage}{CalPolyBeige}{%
      \parbox[c][#1][c]{0.86\linewidth}{%
        \centering
        {\color{DeckMuted}\scriptsize Image placeholder: }%
        {\color{DeckInk}\scriptsize\ttfamily\detokenize{#2}}%
      }%
    }%
  }%
}

\newcommand{\sectionframe}[1]{
  \begingroup
  \setbeamercolor{background canvas}{bg=CalPolyGreen}
  \begin{frame}[plain]
    \vfill
    {\color{white}\LARGE\bfseries #1\par}
    \vspace{0.18in}
    {\color{CalPolyGold}\rule{0.22\paperwidth}{3pt}}
    {\color{white}\rule{0.62\paperwidth}{3pt}}
    \vfill
  \end{frame}
  \endgroup
}

\newcommand{\placeholdernote}{
  \begin{block}{Placeholder}
    Replace these bullets with the final lecture material, examples, and in-class activities.
  \end{block}
}


\pdfstringdefDisableCommands{\def\translate#1{#1}}

\graphicspath{{../assets/lecture-21/}}

% If an image file is not yet available, the slide displays a labeled
% placeholder using the intended filename. Place the matching file in
% ../assets/lecture-21/ to replace the placeholder automatically.
\newcommand{\lectureimage}[2][1.75in]{%
  \IfFileExists{../assets/lecture-21/#2}{%
    \includegraphics[width=\linewidth,height=#1,keepaspectratio]{#2}%
  }{%
    \fbox{%
      \parbox[c][#1][c]{0.94\linewidth}{%
        \centering
        \small\textbf{Image Placeholder}\\[0.08in]
        \scriptsize\texttt{\detokenize{#2}}
      }%
    }%
  }%
}

\setbeamersize{text margin left=0.42in,text margin right=0.42in}

\coursenumber{CS 1001}
\coursetitle{Introduction to Computing}
\lecturenumber{21}
\lecturetitle{Tracing Mutation and Side Effects}
\lecturedate{Fall 2026}
\courseunit{8. References and Mutation}

\begin{document}

\makedecktitle

\begin{frame}{Today}
  \begin{itemize}
    \item Reinforce the distinction between \textbf{aliasing}, \textbf{copying}, and \textbf{rebinding}.
    \item Predict the effects of functions that receive and mutate lists.
    \item Identify when surprising output is caused by shared references and side effects.
    \item Repair unintended mutation by copying or constructing independent data.
    \item Distinguish accidental mutation from intentional, documented mutation.
    \item Simulate changes to immutable strings and tuples by constructing new objects.
  \end{itemize}

  \begin{keyidea}
    The central tracing question is not merely ``what value does this variable have?'' It is also ``which object does this name currently refer to?''
  \end{keyidea}
\end{frame}

\sectionframe{Part 1: A Reference-Aware Tracing Model}

\begin{frame}{Why Some Programs Become Hard to Predict}
  \begin{itemize}
    \item With immutable values, a reassignment typically affects only one binding.
    \item With mutable lists, several names may refer to the same object.
    \item A mutation through one alias changes the object observed through every alias.
    \item Function parameters can temporarily become additional aliases to caller-owned lists.
  \end{itemize}

  \begin{alertblock}{Source of confusion}
    A line of code may mention only one variable name while changing state visible through several different names.
  \end{alertblock}
\end{frame}

\begin{frame}{Three Operations That Must Remain Distinct}
  \begin{center}
    \renewcommand{\arraystretch}{1.3}
    \begin{tabular}{l|l|l}
      \textbf{Operation} & \textbf{Example} & \textbf{Effect} \\\hline
      Rebinding & \texttt{x = new\_value} & \texttt{x} refers to another object \\
      Aliasing & \texttt{y = x} & another name refers to the same object \\
      Mutation & \texttt{x[0] = 7} & existing mutable object changes
    \end{tabular}
  \end{center}

  \begin{keyidea}
    Assignment does not imply copying. Indexed assignment to a list does not rebind the list variable.
  \end{keyidea}
\end{frame}

\begin{frame}{Reminder: Aliasing}
  \begin{columns}[T,onlytextwidth]
    \begin{column}{0.52\textwidth}
      \begin{block}{Code}
        \texttt{a = [10, 20, 30]}\\
        \texttt{b = a}
      \end{block}

      After the second statement:
      \begin{itemize}
        \item there is one list object;
        \item both \texttt{a} and \texttt{b} refer to it;
        \item \texttt{a is b} is \texttt{True}.
      \end{itemize}
    \end{column}
    \begin{column}{0.40\textwidth}
      \lectureimage[2.55in]{alias-reference-diagram.png}
    \end{column}
  \end{columns}
\end{frame}

\begin{frame}{A Reliable Tracing Procedure}
  For each statement involving a list:
  \begin{enumerate}
    \item Identify the object or value produced by the right-hand side.
    \item Determine whether the statement creates a new object, a new alias, or a mutation.
    \item Update the name-to-object bindings.
    \item Update the state of any mutated object.
    \item If a function is called, add its parameter bindings to the same object diagram.
  \end{enumerate}

  \begin{keyidea}
    Track \textbf{bindings} and \textbf{object state} separately.
  \end{keyidea}
\end{frame}

\begin{frame}{Warm-Up: Predict Before Running}
  \begin{block}{Code}
    \texttt{first = [2, 4, 6]}\\
    \texttt{second = first}\\
    \texttt{second[1] = 99}\\
    \texttt{print(first)}\\
    \texttt{print(second)}
  \end{block}

  \begin{activity}
    Predict both printed lists. Draw the reference diagram before deciding.
  \end{activity}
\end{frame}

\begin{frame}{Warm-Up: Reveal}
  \begin{block}{Output}
    \texttt{[2, 99, 6]}\\
    \texttt{[2, 99, 6]}
  \end{block}

  \begin{itemize}
    \item \texttt{second = first} created an alias, not a copy.
    \item \texttt{second[1] = 99} mutated the shared list object.
    \item The expression \texttt{first} observes that same mutated object.
  \end{itemize}
\end{frame}

\sectionframe{Part 2: Functions Can Create Caller-Visible Side Effects}

\begin{frame}{Prediction 1: A Function Mutates Its Parameter}
  \begin{block}{Code}
    \texttt{def zero\_first(values: list[int]) -> None:}\\
    \hspace{0.18in}\texttt{values[0] = 0}\\[0.08in]
    \texttt{scores = [8, 9, 10]}\\
    \texttt{alias = scores}\\
    \texttt{zero\_first(alias)}\\
    \texttt{print(scores)}\\
    \texttt{print(alias)}
  \end{block}

  \begin{activity}
    Predict both outputs. At the indexed assignment, which names refer to the same list object?
  \end{activity}
\end{frame}

\begin{frame}{Prediction 1: What Actually Happens}
  During the function call:
  \begin{center}
    \texttt{scores} $\searrow$\qquad
    \texttt{alias} $\rightarrow$ \boxed{[8, 9, 10]} $\leftarrow$ \texttt{values}
  \end{center}

  After \texttt{values[0] = 0}:
  \begin{center}
    \boxed{[0, 9, 10]}
  \end{center}

  \begin{block}{Output}
    \texttt{[0, 9, 10]}\\
    \texttt{[0, 9, 10]}
  \end{block}

  \begin{keyidea}
    Parameter binding can introduce a temporary alias to a caller-owned mutable object.
  \end{keyidea}
\end{frame}

\begin{frame}{The Unexpected Result Is a Side Effect}
  \begin{cpdefinition}
    A \textbf{side effect} is an observable change to program state that occurs in addition to producing a return value.
  \end{cpdefinition}

  In the previous function:
  \begin{itemize}
    \item the function returns \texttt{None};
    \item the caller's list changes anyway;
    \item the change remains after the call frame ends.
  \end{itemize}

  \begin{alertblock}{Important}
    The problem is not that mutation occurred. The problem is that mutation may be surprising if the function's purpose does not clearly promise it.
  \end{alertblock}
\end{frame}

\begin{frame}{Repair Strategy 1: Copy Before Calling}
  \begin{block}{Code}
    \texttt{scores = [8, 9, 10]}\\
    \texttt{working = scores.copy()}\\
    \texttt{zero\_first(working)}
  \end{block}

  \begin{center}
    \texttt{scores} $\rightarrow$ \boxed{[8, 9, 10]}\\[0.12in]
    \texttt{working} $\rightarrow$ \boxed{[0, 9, 10]}
  \end{center}

  \begin{keyidea}
    Copying breaks the alias before mutation occurs.
  \end{keyidea}
\end{frame}

\begin{frame}{Repair Strategy 2: Copy Inside the Function}
  \begin{block}{Code}
    \texttt{def with\_zero\_first(values: list[int]) -> list[int]:}\\
    \hspace{0.18in}\texttt{result = values.copy()}\\
    \hspace{0.18in}\texttt{result[0] = 0}\\
    \hspace{0.18in}\texttt{return result}
  \end{block}

  \begin{itemize}
    \item The parameter \texttt{values} still aliases the caller's list.
    \item \texttt{result} refers to a distinct list object.
    \item Mutation is applied only to the new list.
  \end{itemize}

  \begin{keyidea}
    A ``return a modified copy'' design makes the state change explicit in the returned value.
  \end{keyidea}
\end{frame}

\sectionframe{Part 3: Harder Aliasing Traces}

\begin{frame}{Prediction 2: A Local Alias Is Still an Alias}
  \begin{block}{Code}
    \texttt{def adjust(values: list[int]) -> list[int]:}\\
    \hspace{0.18in}\texttt{local = values}\\
    \hspace{0.18in}\texttt{local[1] = local[1] + 10}\\
    \hspace{0.18in}\texttt{return local}\\[0.08in]
    \texttt{data = [1, 2, 3]}\\
    \texttt{result = adjust(data)}\\
    \texttt{result[-1] = 99}\\
    \texttt{print(data)}\\
    \texttt{print(result)}
  \end{block}

  \begin{activity}
    Predict both outputs. Is \texttt{local = values} a copy operation?
  \end{activity}
\end{frame}

\begin{frame}{Prediction 2: Reveal and Diagnose}
  \begin{block}{Output}
    \texttt{[1, 12, 99]}\\
    \texttt{[1, 12, 99]}
  \end{block}

  \begin{itemize}
    \item \texttt{local = values} creates another name for the same object.
    \item The function mutates that shared object at index 1.
    \item Returning \texttt{local} returns a reference to the same list.
    \item Therefore \texttt{result} becomes yet another alias.
    \item Mutating \texttt{result[-1]} also changes what \texttt{data} observes.
  \end{itemize}
\end{frame}

\begin{frame}{Fix Prediction 2}
  Replace the aliasing assignment:
  \begin{block}{Problem}
    \texttt{local = values}
  \end{block}

  with an independent copy:
  \begin{block}{Repair}
    \texttt{local = values.copy()}
  \end{block}

  Then:
  \begin{center}
    \texttt{data} $\rightarrow$ \boxed{[1, 2, 3]}\\[0.12in]
    \texttt{result} $\rightarrow$ \boxed{[1, 12, 99]}
  \end{center}
\end{frame}

\begin{frame}{Prediction 3: Concatenation Can Break an Alias}
  \begin{block}{Code}
    \texttt{def add\_marker(values: list[int]) -> list[int]:}\\
    \hspace{0.18in}\texttt{temp = values}\\
    \hspace{0.18in}\texttt{temp = temp + [0]}\\
    \hspace{0.18in}\texttt{temp[0] = 9}\\
    \hspace{0.18in}\texttt{return temp}\\[0.08in]
    \texttt{original = [1, 2, 3]}\\
    \texttt{changed = add\_marker(original)}\\
    \texttt{print(original)}\\
    \texttt{print(changed)}
  \end{block}

  \begin{activity}
    Does \texttt{original} change? Identify the exact statement where the reference relationship changes.
  \end{activity}
\end{frame}

\begin{frame}{Prediction 3: Reveal}
  \begin{block}{Output}
    \texttt{[1, 2, 3]}\\
    \texttt{[9, 2, 3, 0]}
  \end{block}

  \begin{itemize}
    \item Initially, \texttt{temp} aliases \texttt{values}.
    \item \texttt{temp + [0]} constructs a new list.
    \item The assignment then rebinds \texttt{temp} to that new list.
    \item The later indexed mutation affects only the new list.
  \end{itemize}

  \begin{keyidea}
    An alias can cease to be an alias when one name is rebound to a different object.
  \end{keyidea}
\end{frame}

\begin{frame}{Compare: \texttt{+} Versus \texttt{extend}}
  \begin{center}
    \renewcommand{\arraystretch}{1.35}
    \begin{tabular}{l|l|l}
      \textbf{Operation} & \textbf{Creates new list?} & \textbf{Mutates original?} \\\hline
      \texttt{new = values + [0]} & yes & no \\
      \texttt{values.extend([0])} & no & yes
    \end{tabular}
  \end{center}

  \begin{alertblock}{Tracing consequence}
    With aliases, \texttt{extend} is visible through every reference to that list. Concatenation produces a separate list object.
  \end{alertblock}
\end{frame}

\begin{frame}{Prediction 4: Copy, Then Mutate}
  \begin{block}{Code}
    \texttt{def adjust\_copy(values: list[int]) -> list[int]:}\\
    \hspace{0.18in}\texttt{result = values[:]}\\
    \hspace{0.18in}\texttt{result[0] = 50}\\
    \hspace{0.18in}\texttt{return result}\\[0.08in]
    \texttt{a = [5, 6, 7]}\\
    \texttt{b = a}\\
    \texttt{c = adjust\_copy(b)}\\
    \texttt{b[1] = 60}
  \end{block}

  \begin{activity}
    Predict the final values observed through \texttt{a}, \texttt{b}, and \texttt{c}.
  \end{activity}
\end{frame}

\begin{frame}{Prediction 4: Reveal}
  \begin{center}
    \renewcommand{\arraystretch}{1.35}
    \begin{tabular}{l|l|l}
      \textbf{Name} & \textbf{Final value observed} & \textbf{Object} \\\hline
      \texttt{a} & \texttt{[5, 60, 7]} & L1 \\
      \texttt{b} & \texttt{[5, 60, 7]} & L1 \\
      \texttt{c} & \texttt{[50, 6, 7]} & L2
    \end{tabular}
  \end{center}

  \begin{keyidea}
    The copy separates \texttt{c} from the aliased pair \texttt{a}/\texttt{b}; it does not retroactively remove aliasing between \texttt{a} and \texttt{b}.
  \end{keyidea}
\end{frame}

\sectionframe{Part 4: Techniques for Avoiding Unintended Shared Mutation}

\begin{frame}{Several Ways to Create an Independent List}
  For the flat lists of primitive values used in this unit, each expression below creates a distinct list object:
  \begin{block}{Examples}
    \texttt{copy1 = original.copy()}\\
    \texttt{copy2 = original[:]}\\
    \texttt{copy3 = list(original)}
  \end{block}

  For a known fixed structure, explicit recreation is also possible:
  \begin{block}{Explicit recreation}
    \texttt{copy4 = [original[0], original[1], original[2]]}
  \end{block}
\end{frame}

\begin{frame}{Mental Model: Copying Breaks the Shared Reference}
  \begin{columns}[T,onlytextwidth]
    \begin{column}{0.48\textwidth}
      \begin{block}{Before copy}
        \texttt{a} $\searrow$\\
        \hspace{0.55in}\boxed{[3, 4, 5]}\\[-0.02in]
        \texttt{b} $\nearrow$
      \end{block}

      \begin{block}{After \texttt{b = a.copy()}}
        \texttt{a} $\rightarrow$ \boxed{[3, 4, 5]}\\
        \texttt{b} $\rightarrow$ \boxed{[3, 4, 5]}
      \end{block}
    \end{column}
    \begin{column}{0.44\textwidth}
      \lectureimage[2.7in]{copy-breaks-alias.png}
    \end{column}
  \end{columns}
\end{frame}

\begin{frame}{Technique 1: \texttt{list.copy()}}
  \begin{block}{Code}
    \texttt{original = [10, 20, 30]}\\
    \texttt{working = original.copy()}\\
    \texttt{working[0] = 99}
  \end{block}

  \begin{block}{Result}
    \texttt{original == [10, 20, 30]}\\
    \texttt{working == [99, 20, 30]}
  \end{block}

  \begin{keyidea}
    \texttt{copy()} directly communicates the programmer's intent to create another list object.
  \end{keyidea}
\end{frame}

\begin{frame}{Technique 2: Full Slicing}
  \begin{block}{Code}
    \texttt{original = [10, 20, 30]}\\
    \texttt{working = original[:]}\\
    \texttt{working[-1] = 0}
  \end{block}

  A full slice selects all elements and constructs a new list.

  \begin{block}{Result}
    \texttt{original == [10, 20, 30]}\\
    \texttt{working == [10, 20, 0]}
  \end{block}
\end{frame}

\begin{frame}{Technique 3: \texttt{list(original)}}
  \begin{block}{Code}
    \texttt{original = [1, 2, 3]}\\
    \texttt{working = list(original)}\\
    \texttt{working.remove(2)}
  \end{block}

  \begin{itemize}
    \item \texttt{list(original)} constructs a new list from the supplied sequence.
    \item The new list has equal contents initially but distinct identity.
  \end{itemize}

  \begin{block}{Identity check}
    \texttt{working == original} is initially \texttt{True}.\\
    \texttt{working is original} is \texttt{False}.
  \end{block}
\end{frame}

\begin{frame}{Technique 4: Explicit Reconstruction}
  When the structure is small and fixed, the new value can be written explicitly.

  \begin{block}{Code}
    \texttt{point = [4, 7, 9]}\\
    \texttt{changed = [point[0], 100, point[2]]}
  \end{block}

  \begin{itemize}
    \item \texttt{point} is unchanged.
    \item \texttt{changed} is a new list with the desired contents.
    \item The intended change is represented as construction, not mutation.
  \end{itemize}
\end{frame}

\begin{frame}{Where Should the Copy Happen?}
  There are two legitimate designs:
  \begin{columns}[T,onlytextwidth]
    \begin{column}{0.47\textwidth}
      \begin{block}{Caller protects itself}
        \texttt{working = data.copy()}\\
        \texttt{mutate(working)}
      \end{block}
      Useful when the called function is intentionally mutating.
    \end{column}
    \begin{column}{0.47\textwidth}
      \begin{block}{Function protects caller}
        \texttt{def transformed(data):}\\
        \hspace{0.14in}\texttt{result = data.copy()}\\
        \hspace{0.14in}\texttt{...}\\
        \hspace{0.14in}\texttt{return result}
      \end{block}
      Useful when the function promises not to modify its argument.
    \end{column}
  \end{columns}
\end{frame}

\sectionframe{Part 5: Mutation Is Not Inherently Bad}

\begin{frame}{Intentional Mutation Is a Valid Design Choice}
  \begin{keyidea}
    Mutation is appropriate when changing an existing object is the intended behavior of the operation.
  \end{keyidea}

  Examples already used in this unit:
  \begin{itemize}
    \item \texttt{append} intentionally grows a list;
    \item \texttt{remove} intentionally deletes an element;
    \item \texttt{reverse} intentionally changes list order;
    \item indexed assignment intentionally replaces an element.
  \end{itemize}

  The important question is whether that state change is part of the function's clear contract.
\end{frame}

\begin{frame}{Intentional Example: Swap the Endpoints}
  \begin{block}{Code}
    \texttt{def swap\_ends(values: list[int]) -> None:}\\
    \hspace{0.18in}\texttt{first = values[0]}\\
    \hspace{0.18in}\texttt{values[0] = values[-1]}\\
    \hspace{0.18in}\texttt{values[-1] = first}
  \end{block}

  \begin{block}{Use}
    \texttt{numbers = [10, 20, 30, 40]}\\
    \texttt{swap\_ends(numbers)}
  \end{block}

  After the call, \texttt{numbers} is intentionally \texttt{[40, 20, 30, 10]}.
\end{frame}

\begin{frame}{Why a \texttt{None} Return Can Be Appropriate}
  For a function whose primary purpose is mutation:
  \begin{itemize}
    \item the important observable result is the changed object;
    \item returning \texttt{None} emphasizes that the function acts through a side effect;
    \item the caller should not expect a separate transformed list.
  \end{itemize}

  \begin{block}{Contrast}
    \texttt{swap\_ends(values) -> None} \hfill mutates supplied list\\
    \texttt{swapped\_ends(values) -> list[int]} \hfill returns a new list
  \end{block}
\end{frame}

\begin{frame}{Two Clear Interface Designs}
  \begin{center}
    \renewcommand{\arraystretch}{1.35}
    \begin{tabular}{l|l|l}
      \textbf{Design} & \textbf{Argument} & \textbf{Returned value} \\\hline
      Mutating operation & may change & often \texttt{None} \\
      Transform-and-return & unchanged & new value/object
    \end{tabular}
  \end{center}

  \begin{alertblock}{Avoid ambiguity}
    A function is easier to reason about when its name, purpose statement, tests, and implementation agree about whether mutation is expected.
  \end{alertblock}
\end{frame}

\sectionframe{Part 6: Simulated Mutation for Immutable Sequences}

\begin{frame}{Strings Cannot Be Mutated In Place}
  \begin{block}{Invalid code}
    \texttt{text = "cat"}\\
    \texttt{text[0] = "b"}
  \end{block}

  Python raises a \texttt{TypeError} because strings are immutable.

  \begin{keyidea}
    To represent a changed string, construct a new string and bind a name to it.
  \end{keyidea}
\end{frame}

\begin{frame}{Simulated String Mutation}
  \begin{block}{Construct a new string}
    \texttt{text = "cat"}\\
    \texttt{changed = "b" + text[1:]}
  \end{block}

  Now:
  \begin{center}
    \texttt{text} $\rightarrow$ \boxed{\texttt{"cat"}}\\[0.10in]
    \texttt{changed} $\rightarrow$ \boxed{\texttt{"bat"}}
  \end{center}

  Reusing the same variable name is also possible:
  \begin{block}{Rebinding}
    \texttt{text = "b" + text[1:]}
  \end{block}

  This is construction plus rebinding, not mutation of the original string.
\end{frame}

\begin{frame}{Tuples Are Immutable Sequences}
  \begin{block}{Example}
    \texttt{point = (10, 20, 30)}
  \end{block}

  Tuples support familiar sequence operations:
  \begin{itemize}
    \item indexing: \texttt{point[0]};
    \item negative indexing: \texttt{point[-1]};
    \item slicing: \texttt{point[1:]};
    \item length and membership checks.
  \end{itemize}

  But indexed assignment is not permitted:
  \begin{block}{Invalid}
    \texttt{point[1] = 99}
  \end{block}
\end{frame}

\begin{frame}{Simulated Tuple Mutation}
  Suppose we want to replace the middle element:
  \begin{block}{Construct a new tuple}
    \texttt{point = (10, 20, 30)}\\
    \texttt{changed = point[:1] + (99,) + point[2:]}
  \end{block}

  \begin{itemize}
    \item \texttt{point} remains \texttt{(10, 20, 30)}.
    \item \texttt{changed} is \texttt{(10, 99, 30)}.
    \item \texttt{(99,)} is a one-element tuple; the comma is required.
  \end{itemize}
\end{frame}

\begin{frame}{Returning a Changed Tuple}
  \begin{block}{Code}
    \texttt{def replace\_middle(values: tuple[int, int, int],}\\
    \hspace{1.22in}\texttt{new\_value: int) -> tuple[int, int, int]:}\\
    \hspace{0.18in}\texttt{return values[:1] + (new\_value,) + values[2:]}
  \end{block}

  \begin{itemize}
    \item The input tuple is never mutated.
    \item The function returns a different tuple containing the requested change.
    \item This is the immutable analogue of ``copy, change, return.''
  \end{itemize}
\end{frame}

\sectionframe{Part 7: Integrated Tracing Challenge}

\begin{frame}{Integrated Prediction}
  \begin{block}{Code}
    \texttt{def bump\_first(values: list[int]) -> None:}\\
    \hspace{0.18in}\texttt{values[0] = values[0] + 1}\\[0.06in]
    \texttt{def adjusted(values: list[int]) -> list[int]:}\\
    \hspace{0.18in}\texttt{result = values[:]}\\
    \hspace{0.18in}\texttt{result[-1] = 0}\\
    \hspace{0.18in}\texttt{bump\_first(result)}\\
    \hspace{0.18in}\texttt{return result}\\[0.08in]
    \texttt{a = [5, 6, 7]}\\
    \texttt{b = a}\\
    \texttt{c = adjusted(b)}\\
    \texttt{b[1] = 9}
  \end{block}

  \begin{activity}
    Predict the final values of \texttt{a}, \texttt{b}, and \texttt{c}. Draw object labels before tracing statements.
  \end{activity}
\end{frame}

\begin{frame}{Integrated Trace: Object Structure}
  After \texttt{b = a}:
  \begin{center}
    \texttt{a} $\searrow$\qquad \boxed{L1 = [5, 6, 7]} \qquad$\swarrow$ \texttt{b}
  \end{center}

  Inside \texttt{adjusted}:
  \begin{itemize}
    \item parameter \texttt{values} aliases L1;
    \item \texttt{result = values[:]} creates L2;
    \item \texttt{result[-1] = 0} changes only L2;
    \item \texttt{bump\_first(result)} temporarily aliases L2 through its parameter and mutates L2.
  \end{itemize}

  \begin{center}
    L1 = \texttt{[5, 6, 7]} \qquad L2 = \texttt{[6, 6, 0]}
  \end{center}
\end{frame}

\begin{frame}{Integrated Trace: Final Result}
  After \texttt{c = adjusted(b)}:
  \begin{center}
    \texttt{a, b} $\rightarrow$ L1 = \texttt{[5, 6, 7]}\\[0.10in]
    \texttt{c} $\rightarrow$ L2 = \texttt{[6, 6, 0]}
  \end{center}

  Then \texttt{b[1] = 9} mutates L1.

  \begin{block}{Final values}
    \texttt{a == [5, 9, 7]}\\
    \texttt{b == [5, 9, 7]}\\
    \texttt{c == [6, 6, 0]}
  \end{block}
\end{frame}

\begin{frame}{Debugging Unexpected Mutation}
  When a list changes unexpectedly, ask these questions in order:
  \begin{enumerate}
    \item Which statement first changed the object's state?
    \item Which names referred to that object at that moment?
    \item Did a function parameter alias a caller-owned list?
    \item Was the mutation intentional according to the function's contract?
    \item If not, where should an independent object have been created?
  \end{enumerate}

  \begin{keyidea}
    Locate the first unintended mutation, not merely the final surprising print statement.
  \end{keyidea}
\end{frame}

\begin{frame}{Common Incorrect Mental Models}
  \begin{itemize}
    \item \textbf{Incorrect:} ``Assignment copies the list.''\\
      \textbf{Correct:} ordinary assignment copies the reference/binding relationship.
    \item \textbf{Incorrect:} ``The function changed the variable in the caller.''\\
      \textbf{Correct:} the function mutated an object that both caller and parameter referenced.
    \item \textbf{Incorrect:} ``Returning the list makes a copy.''\\
      \textbf{Correct:} returning an existing list returns a reference to that object.
    \item \textbf{Incorrect:} ``Every operation that looks like a change is mutation.''\\
      \textbf{Correct:} strings and tuples require construction of new objects.
  \end{itemize}
\end{frame}

\begin{frame}{Takeaways}
  \begin{itemize}
    \item Aliasing means multiple names refer to one object.
    \item Mutation changes an existing object and is visible through all of its aliases.
    \item Function parameters can alias mutable objects owned by the caller.
    \item Unexpected shared mutation is a side effect worth diagnosing explicitly.
    \item Use \texttt{copy()}, full slicing, \texttt{list(...)}, or explicit reconstruction when independent list state is required.
    \item Intentional mutation is appropriate when it is part of the function's contract.
    \item Strings and tuples are immutable; represent changes by constructing and rebinding new values.
  \end{itemize}

  \begin{keyidea}
    When tracing mutable data, follow the objects first and the variable names second.
  \end{keyidea}
\end{frame}

\end{document}
